\subsection{TANGENT LAWS OF COMPOSITION}

Let X and Y be manifolds of class \(C^{r}\) . We know (Differentiable and Analytic Manifolds, R, 8.1.4) that \(X \times Y\) is a manifold of class \(C^{r}\) and that the mapping \((\mathrm{T}(\mathrm{pr}_{1}), \mathrm{T}(\mathrm{pr}_{2}))\) , the product of the tangent mappings to the canonical projections, is an isomorphism of class \(C^{r-1}\) of \(\mathrm{T}(\mathrm{X} \times \mathrm{Y})\) onto \(\mathrm{T}(\mathrm{X}) \times \mathrm{T}(\mathrm{Y})\) . This isomorphism is compatible with the vector bundle structures with base space \(X \times Y\) and allows us to identify \(\mathrm{T}(\mathrm{X} \times \mathrm{Y})\) with \(\mathrm{T}(\mathrm{X}) \times \mathrm{T}(\mathrm{Y})\) . Let \(a \in \mathbf{X}, b \in \mathbf{Y}, u \in \mathrm{T}_a(\mathbf{X}), v \in \mathrm{T}_b(\mathbf{Y})\) ; the above identification allows us to consider \((u, v)\) as an element of \(\mathrm{T}_{(a,b)}(\mathbf{X} \times \mathbf{Y})\) ; then

\[
(u, v) = (u, 0) + (0, v)
\]

and \((u,0)\) (resp. \((0,v))\) is the image of \(u\) (resp. \(v\) ) under the tangent mapping to the immersion \(x\mapsto (x,b)\) (resp. \(y\mapsto (a,y)\) ) of X (resp. Y) into \(\mathbf{X}\times \mathbf{Y}\) . When it is necessary to be precise, we shall write \(0_{a}\) for the zero element of \(\mathrm{T}_a(X)\) .

Now let X, Y, Z be manifolds of class \(C^{r}\) and f a mapping of class \(C^{r}\) of \(X \times Y\) into Z. The tangent mapping is, using the above identification, a mapping of class \(C^{r-1}\) of \(T(X) \times T(Y)\) into \(T(Z)\) . For \(u \in T_{a}(X)\) and \(v \in T_{b}(Y)\) ,

\begin{enumerate}
\def\labelenumi{(\arabic{enumi})}
\tightlist
\item
\end{enumerate}

\[
\mathrm{T} (f) (u, v) = \mathrm{T} (f) (u, 0 _ {b}) + \mathrm{T} (f) (0 _ {a}, v),
\]

\[
\mathrm{T} (f) \left(0 _ {a}, 0 _ {b}\right) = 0 _ {f (a, b)}.\tag{2}
\]

On the other hand, the mapping \(y \mapsto f(a, y)\) is the composition of the immersion \(y \mapsto (a, y)\) and f; it follows that

\begin{enumerate}
\def\labelenumi{(\arabic{enumi})}
\setcounter{enumi}{2}
\item
  \(\mathrm{T}(f)(0,v)\) is the image of v under the tangent mapping to \(y\mapsto f(a,y)\) . Similarly
\item
  \(\mathbf{T}(f)(u,0)\) is the image of \(u\) under the tangent mapping to \(x\mapsto f(x,b)\) .
\end{enumerate}

If the mapping \(f\) of \(\mathbf{X} \times \mathbf{Y}\) into \(\mathbf{Z}\) is denoted by \((x, y) \mapsto xy\) , \(uv\) is often used to denote the element \(T(f)(u, v)\) for \(u \in T(\mathbf{X})\) , \(v \in T(\mathbf{Y})\) .

Let X be a manifold of class \(C^{r}\) and \(m: X \times X \to X\) a law of composition of class \(C^{r}\) on X. Then \(T(m)\) is a law of composition of class \(C^{r-1}\) on \(T(X)\) . It is called the law of composition tangent to m. The canonical projection p of \(T(X)\) onto X is compatible with the laws m and \(T(m)\) ; in other words,

\[
p \circ \mathrm{T} (m) = m \circ (p \times p).\tag{5}
\]

It follows from (2) that

\[
\mathrm{T} (m) \left(0 _ {x}, 0 _ {y}\right) = 0 _ {m (x, y)}\tag{6}
\]

for all \(x, y\) in \(\mathbf{X}\) ; in other words, the zero section \(x \mapsto 0_x\) of \(\mathrm{T}(\mathbf{X})\) is compatible with the laws \(m\) and \(\mathrm{T}(m)\) .

PROPOSITION 1. Let X be a manifold of class \(\mathbf{C}^r\) and \(m\) a law of composition of class \(\mathbf{C}^r\) on X. If \(m\) is associative (resp. commutative), then \(\mathrm{T}(m)\) is associative (resp. commutative).

If \(m\) is associative, then \(m \circ (m \times \mathrm{Id}_{\mathbf{x}}) = m \circ (\mathrm{Id}_{\mathbf{x}} \times m)\) , whence

\[
\mathrm{T} (m) \circ (\mathrm{T} (m) \times \mathrm{Id} _ {\mathrm{T} (\mathbf {X})}) = \mathrm{T} (m) \circ (\mathrm{Id} _ {\mathrm{T} (\mathbf {X})} \times \mathrm{T} (m))
\]

and hence \(\mathrm{T}(m)\) is associative. Let \(s\) be the mapping \((x,y)\mapsto (y,x)\) of \(\mathbf{X}\times \mathbf{X}\) into \(\mathbf{X}\times \mathbf{X}\) . If \(m\) is commutative, then \(m\circ s = m\) and hence

\[
\mathrm{T} (m) \circ \mathrm{T} (s) = \mathrm{T} (m).
\]

But \(\mathrm{T}(s)\) is the mapping \((u,v)\mapsto (v,u)\) of \(\mathrm{T}(\mathbf{X})\times \mathrm{T}(\mathbf{X})\) into \(\mathrm{T}(\mathbf{X})\times \mathrm{T}(\mathbf{X})\) . Hence \(\mathrm{T}(m)\) is commutative.

PROPOSITION 2. Let X be a manifold of class \(\mathbf{C}^r\) , \(m\) a law of composition of class \(\mathbf{C}^r\) on X and \(e\) an identity element for \(m\) .

\begin{enumerate}
\def\labelenumi{(\roman{enumi})}
\item
  The vector \(0_{e}\) is an identity element for \(\mathbf{T}(m)\) .
\item
  \(\mathrm{T}_{e}(\mathbf{X})\) is stable under \(\mathrm{T}(m)\) and the law of composition induced on \(\mathrm{T}_{e}(\mathbf{X})\) by \(\mathrm{T}(m)\) is the vector space addition on \(\mathrm{T}_{e}(\mathbf{X})\) .
\item
  Let U be an open subset of X and \(\alpha\) a mapping of class \(C^{r}\) of U into X such that, for all \(x \in U\) , \(\alpha(x)\) is the inverse of x under m. Then, for all \(u \in T(U)\) , \(T(\alpha)u\) is the inverse of u under \(T(m)\) .
\end{enumerate}

Properties (3) and (4) show that \(\mathrm{T}(m)(0_e,u) = \mathrm{T}(m)(u,0_e) = u\) for all \(u\in \mathrm{T}(\mathbf{X})\) , whence (i). For \(u,v\) in \(\mathrm{T}_e(\mathbf{X})\) ,

\[
\mathrm{T} (m) (u, v) = \mathrm{T} (m) (u, 0 _ {e}) + \mathrm{T} (m) (0 _ {e}, v) = u + v,
\]

whence (ii). Finally the relations \(m(x, \alpha(x)) = m(\alpha(x), x) = e\) for all \(x \in \mathbf{U}\) imply

\[
\mathrm{T} (m) (u, \mathrm{T} (\alpha) (u)) = \mathrm{T} (m) (\mathrm{T} (\alpha) u, u) = 0 _ {e}
\]

for all \(u \in \mathrm{T}(\mathrm{U})\) , whence (iii).

PROPOSITION 3. Let \(\mathbf{X}_1, \mathbf{X}_2, \ldots, \mathbf{X}_p\) , Y be manifolds of class \(\mathbf{C}^r\) , i an integer of \([1, p]\) , \(m_t\) (resp. n) a law of composition of class \(\mathbf{C}^r\) on \(\mathbf{X}_t\) (resp. Y) and u a mapping of class \(\mathbf{C}^r\) of \(\mathbf{X}_1 \times \mathbf{X}_2 \times \dots \times \mathbf{X}_p\) into Y. If u is distributive relative to the variable of index i, then T(u) is distributive relative to the variable of index i.

The proof is analogous to that of Proposition 1.

